\section{What the Approval Obliges---and What It Does Not}

\subsection{Public Revelation and private revelation}

The doctrinal frame is fixed and prior to every fact in this
study. Public Revelation is complete in Christ: ``no new public
revelation is to be expected before the glorious manifestation of
our Lord Jesus Christ.'' Private revelations, even recognized
ones, ``do not belong\ldots\ to the deposit of faith. It is not
their role to improve or complete Christ's definitive Revelation,
but to help live more fully by it in a certain period of
history.''\endnote{\work{Catechism of the Catholic Church}, 2nd
ed., nn.~66--67 (n.~66 quoting \work{Dei Verbum} 4), English
text read 2026-07-25 at
\url{https://www.vatican.va/archive/ENG0015/__PH.HTM} (the
Vatican web edition of CCC 65--67).} Lourdes falls entirely on
the private-revelation side of that line. The 1862 judgment
declares the faithful \emph{justified} in believing the
apparition certain; it does not, and could not, make that belief
an article of divine and Catholic faith. The assent invited is
what the tradition calls \term{fides humana}: a prudent human
assent, resting on the quality of the evidence and the Church's
discernment, revisable in principle as human judgments are, and
never the theological faith owed to God revealing.

The Holy See's standing formula for what approval means was
quoted at \S1: the message contains nothing contrary to faith or
morals; it is lawful to publish it; the faithful are authorized
to accept it with prudence. Three verbs---contains, is lawful,
are authorized---and none of them is ``must.'' A Catholic who,
without contempt for the Church's judgment, remains personally
unpersuaded by the Lourdes evidence sins against nothing; a
Catholic who finds in Lourdes a support of faith stands
authorized and, since 1862, on the express judgment of the
competent bishop that the ground is solid. What is not open to
either is to treat the apparition as a second deposit of faith,
or its approval as a dogma.\endnote{Doctrinal synthesis over CCC
66--67, the 2000 CDF explanation as cited in the 2024 norms'
presentation (\S1), and the 1862 operative wording (\term{sont
fond\'es \`a la croire certaine}---a warrant, not a precept).}

\subsection{The exact objects of the standing acts}

Because reliance risk lives in the details, the objects of the
several acts are restated here at their exact scope.

The \emph{1862 judgment} covers the eighteen apparitions of 1858
to Bernadette Soubirous and authorizes diocesan cult (\S4.3). It
does not certify any particular narrative of the events beyond
its own summary, any relic, image, book, or later devotion (its
own art.~2 keeps those under separate approval), any physical
theory of the spring, or any cure beyond the seven it recognized.

The \emph{cure recognitions} (\S5) each cover one named person's
cure and nothing else. They do not make the water curative,
predict anyone else's healing, or convert the Bureau's negative
medical finding into a scientific doctrine. No one is obliged to
believe any particular recognized cure miraculous; the
recognitions authorize thanksgiving, they do not legislate it.

The \emph{liturgical acts} (\S7.1--7.2) establish a memorial and
a world day of prayer; they commit the Church to the celebration,
not to a chronicle. The \emph{canonization} (\S6) infallibly
proposes Bernadette as a saint to be venerated; it leaves the
apparitions where the 1862 act put them, as the current norms'
own principle confirms (\S6.4). And \emph{papal words} at and
about Lourdes---encyclical, homilies, messages---teach with the
authority proper to each genre; none of them is an event-decree.

\subsection{Devotional practice within the boundary}

Everything the approved devotion asks can be stated without
strain inside ordinary Catholic doctrine, and the guardians of
Lourdes have repeatedly said so themselves. The water is
ordinary; its use is a gesture of faith and penance, not a
sacrament, and ``will have no virtue without faith'' (\S2.6). The
baths, the processions, the candles are pious exercises framing
the shrine's actual center, which its own practice makes
unmistakable: the Eucharist and the sacraments. The Eucharistic
procession and blessing of the sick stand at the heart of the
Lourdes day; the cures themselves cluster around the sacraments
as much as the baths---Traynor after the Eucharistic procession,
Moriau in Eucharistic adoration after receiving the Anointing of
the Sick (\S5.4). Pius XII's rule of interpretation governs:
``Everything about Mary directs us to her Son, our only
Savior''; Mary's role at Lourdes, as everywhere, is created,
maternal, subordinate, and wholly derived from Christ's unique
mediation---which is exactly how Benedict XVI preached the
``smile of Mary'' to the sick, as ``a true reflection of God's
tenderness,'' not an independent source of
grace.\endnote{\work{Le p\`elerinage de Lourdes}, n.~23; Benedict
XVI, homily of 15 September 2008; doctrinally, \work{Lumen
gentium} 60--62 (Mary's subordinate mediation) and CCC 970
govern, and nothing in the approved Lourdes corpus strains
them---a fact worth stating because it is not true of every
reported apparition.}

Two pastoral boundaries deserve their own sentences, because
Lourdes' whole modern credibility hangs on them. First: nothing
in the approved devotion may displace medicine. The shrine's own
apparatus is saturated with physicians; declaring a cure begins
with medical documentation; and no organ of Lourdes counsels
anyone to abandon treatment, diagnosis, or public-health
instructions in favor of the baths. Second: the dignity of the
sick person does not depend on being cured. The overwhelming
majority of Lourdes pilgrims who are sick go home sick; the
shrine's theology---from the 1858 promise ``not in this world''
through the World Day of the Sick to Benedict XVI's ``that
dignity which never abandons the sick person''---locates the
grace of Lourdes in conversion, endurance, communion, and hope,
with bodily cure as a rare and unowed sign. Any use of Lourdes
to tell an uncured person that their faith failed inverts the
approved message.\endnote{Editorial synthesis; components at
\S5.5, \S7.2, \S7.4. The 1992 institution letter's purpose
clause---offering one's suffering, seeing in the sick brother
``the face of Christ''---is the controlling pastoral text.}

\subsection{Disagreement preserved}

A judgment-first study must also say plainly where the record
invites, and the Church tolerates, disagreement. The 1862 act
itself argues from evidence---testimony, fruits, cures---and
evidence-based judgments can be weighed differently by serious
minds; the Church has never anathematized the historian who
weighs them otherwise. The nineteenth-century controversies
(rationalist re-readings of Bernadette's psychology, Zola's
literary appropriation, later source-critical disputes about the
harmonized narratives) are part of Lourdes' reception history;
this study has not adjudicated them, and its own corpus
boundaries (Appendix~A) mark exactly where its warrant stops.
What the Church's acts secure is narrower and firmer than
victory in those debates: that a Catholic may, with the Church,
prudently believe; that the devotion is doctrinally safe; and
that the sick are to be honored. Nothing in this study should be
read as claiming more.\endnote{The named controversies are
identified as reception history only; none of their texts was
consulted for this study, and no characterization of their
arguments is offered or should be inferred.}
